\section{LLM Post-Training}\label{sec:llm-posttrain}

Stage~(ii) adapts the pre-trained backbone to EDA scripting, and it is the stage that carries BaiZe's domain claim: the pre-training corpus of Section~\ref{sec:llm-data} is general web, code, and mathematics text and contains no EDA artifacts, so the domain prior must be introduced here.
% 中文：第 (ii) 阶段把预训练骨干适配到 EDA 脚本编写，而承载本文领域主张的正是这一阶段：第 \ref{sec:llm-data} 节的预训练语料是通用网页、代码与数学文本，不含任何 EDA 产物，因此领域先验必须在此引入。
We organise the stage as three steps---\emph{domain annealing}, \emph{supervised fine-tuning}, and \emph{reinforcement learning}---all evaluated under a single frozen protocol (Section~\ref{sec:posttrain-eval}).
% 中文：我们按三个步骤组织该阶段——\emph{领域退火}、\emph{监督微调} 与 \emph{强化学习}——三者在同一套冻结协议下评测（第 \ref{sec:posttrain-eval} 节）。

\subsection{Domain Annealing}\label{sec:posttrain-anneal}
The cheapest lever on the domain prior is to continue pre-training on EDA text before any behavioural tuning.
% 中文：领域退火。对领域先验最廉价的杠杆，是在任何行为层面的微调之前，先用 EDA 文本做继续预训练。
We reuse the warmup--stable--decay machinery of Stage~(i) and anneal on a corpus assembled from PyAether and SKILL API references, EDA tool documentation, and open HDL sources, blended with the general mixture at a controlled ratio.
% 中文：我们复用第 (i) 阶段的预热—稳定—衰减机制，并在由 PyAether 与 SKILL API 参考、EDA 工具文档与开源 HDL 源码汇编而成的语料上做退火，按受控比例与通用混合数据混合。
Because the step changes only the data distribution, we contrast it against an SFT-only control to isolate its contribution from the behavioural tuning that follows.
% 中文：由于该步骤只改变数据分布，我们将其与「仅做 SFT」的对照组进行比较，以把它的贡献从随后的行为微调中隔离出来。

\subsection{Supervised Fine-Tuning}\label{sec:posttrain-sft}
SFT teaches the generate--execute--refine loop that the deployment target requires.
% 中文：监督微调。SFT 教授部署目标所要求的「生成—执行—改进」闭环。
The corpus combines the agentic instruction set (\texttt{UltraData-SFT-Agent-2609}, $51$\,G across 50 shards) and a general instruction set (\texttt{UltraData-SFT-2605}, $319$\,G across 1{,}504 JSONL files, verified byte-identical to the HuggingFace manifest) with PyAether-specific material synthesised from API documentation.
% 中文：语料由智能体指令集（\texttt{UltraData-SFT-Agent-2609}，$51$ G、50 个 shard）与通用指令集（\texttt{UltraData-SFT-2605}，$319$ G、1,504 个 JSONL 文件，已核验与 HuggingFace manifest 逐字节一致）合并，并加入依据 API 文档合成的 PyAether 专属材料。
We ablate the data mixture (agentic / code / mathematics / PyAether), the presence of interleaved tool-call traces, full fine-tuning versus LoRA, and the sequence length: agentic traces---code, tool returns, and successive refinements---routinely exceed the $4096$-token window of Table~\ref{tab:config}, so window extension is treated as a first-class variable rather than a fixed setting.
% 中文：我们消融以下方面：数据配比（智能体／代码／数学／PyAether）、是否包含交错的工具调用轨迹、全量微调还是 LoRA、以及序列长度：智能体轨迹——代码、工具返回与多次迭代改进——通常远超表 \ref{tab:config} 的 $4096$-token 窗口，因此窗口扩展在这里被当作一等变量处理，而非固定设定。

\subsection{Reinforcement Learning with Verifiable Rewards}\label{sec:posttrain-rl}
An execution sandbox turns correctness into a verifiable signal, so we optimise with group-relative policy optimisation and no value network.
% 中文：基于可验证奖励的强化学习。执行沙盒把正确性转化为可验证信号，因此我们采用无需价值网络的组相对策略优化（GRPO）。
Table~\ref{tab:reward} separates the reward variants we study.
% 中文：表 \ref{tab:reward} 区分了我们研究的各种奖励方案。
A binary signal is the simplest baseline; partial credit over executable attempts addresses the sparsity of pass/fail; a tool-call penalty aligns the policy with the harness cost model; and an API-legality term penalises hallucinated symbols.
% 中文：二值信号是最简单的基线；对多次可执行尝试给予部分得分可缓解「通过／不通过」的稀疏性；工具调用惩罚项使策略与 harness 的成本模型对齐；API 合法性项则惩罚虚构的符号。
Because the reward is computed from assertions a policy could in principle target directly, we keep the harness anti-cheat hook enabled in every run and report a cheat rate beside every accuracy figure.
% 中文：由于奖励由策略原则上可以直接瞄准的断言计算而来，我们在每次运行中都保持 harness 的反作弊钩子开启，并在每个准确率数字旁报告作弊率。

\begin{table}[htbp]
\centering
\caption{Reward variants studied in stage~(ii).}
\label{tab:reward}
\begin{tabular}{ll}
\toprule
\textbf{ID} & \textbf{Reward} \\
\midrule
R1 & binary pass/fail \\
R2 & executable $+$ assertion pass ratio \\
R3 & R2 $+$ tool-call penalty \\
R4 & R3 $+$ API-legality term \\
\bottomrule
\end{tabular}
\end{table}

\subsection{Evaluation Protocol}\label{sec:posttrain-eval}

We define four capability milestones to track progress through Stage~(ii):
% 中文：我们定义四个能力里程碑以追踪第 (ii) 阶段的进展：
\begin{itemize}
\item \textbf{M1 (tool-call legality):} the model generates syntactically valid PyAether API calls with no hallucinated symbols (measured by API-legality rate).
\item \textbf{M2 (single-task pass):} the model solves individual EDA scripting tasks end-to-end (Pass@1 on EDA-Eval-PyAether).
\item \textbf{M3 (multi-turn refinement):} the model recovers from execution failures by reading error output and revising code (multi-turn Pass@1 with up to 3 refinement rounds).
\item \textbf{M4 (full benchmark):} the model achieves competitive end-to-end accuracy on the full $158$-task benchmark under the standard five-run protocol.
\end{itemize}
All post-training variants are scored on \textbf{EDA-Eval-PyAether}~\cite{liu2026zhulong}, a benchmark of $158$ real-world EDA scripting tasks.
% 中文：所有后训练变体均在 \textbf{EDA-Eval-PyAether} \cite{liu2026zhulong} 上评分——这是一个包含 $158$ 个真实 EDA 脚本任务的基准。
Each task pairs a natural-language prompt and a function signature with assertion-based test code, and is scored by execution in the vendor sandbox: a task passes when the generated code runs without error and every assertion holds.
% 中文：每个任务把自然语言提示与函数签名同基于断言的测试代码配对，并通过在厂商沙盒中执行来评分：只有当生成的代码无错运行且所有断言成立时，该任务才判为通过。
We follow the benchmark's own protocol so that our numbers are comparable to published ones: five independent runs with reshuffled task order and fresh context, reporting mean Pass@1 and its run-to-run standard deviation, with a $2500$-second per-task timeout and a single non-restartable trace per task.
% 中文：我们遵循该基准自身的协议，以使数字可与已发表结果比较：五次独立运行，打乱任务顺序并重置上下文，报告平均 Pass@1 及其运行间标准差，单任务超时 $2500$ 秒，且每个任务只有一条不可重启的轨迹。

Two controls make the comparison interpretable.
% 中文：两项对照使比较可解释。
First, \emph{contamination}: no task, assertion, or reference solution from the benchmark enters any training split.
% 中文：第一，\emph{污染}：基准中的任何任务、断言或参考解都不得进入任何训练划分。
A blacklist of $193{,}295$ character-level 13-gram fingerprints (FNV-1a 64-bit hashed, NFKC-normalised, with 8-gram fallback for short texts) is built from all $536$ tasks across six evaluation snapshots of EDA-Eval-PyAether; every training shard---including SFT corpora---is scanned for ${\geq}0.80$ n-gram overlap before release, with MinHash($k{=}128$) tracing for source attribution.
% 中文：我们由 EDA-Eval-PyAether 六个评测快照的全部 $536$ 个任务构建了一个包含 $193{,}295$ 条字符级 13-gram 指纹的黑名单（FNV-1a 64 位哈希、NFKC 归一化，短文本回退到 8-gram）；每个训练 shard——包括 SFT 语料——在发布前都被扫描以检查 ${\geq}0.80$ 的 n-gram 重叠，并用 MinHash($k{=}128$) 做来源追溯。
PyAether training material is drawn only from API documentation, and we isolate it by symbol so that benchmark targets cannot leak.
% 中文：PyAether 训练材料仅取自 API 文档，并按符号加以隔离，使基准目标无法泄漏。
Second, \emph{harness parity}: only the backbone is varied. The retrieval stack, sandbox, anti-cheat hook, and agent scaffold are held byte-identical to the published configuration, because the harness supplies most of the observed accuracy and a variant harness would confound the model's contribution with the scaffold's.
% 中文：第二，\emph{harness 一致性}：只改变骨干模型。检索栈、沙盒、反作弊钩子与智能体脚手架均与已发布的配置逐字节保持一致，因为 harness 贡献了观测到的准确率的大部分，若换用变体 harness，就会把模型本身的贡献与脚手架的贡献混在一起。

Results are reported as \textbf{[TBD]}.
% 中文：结果以 \textbf{[TBD]} 表示（待填）。